% !TEX root = ../main.tex
\documentclass[../main.tex]{subfiles}

\begin{document}

\intermezzo{Social Media: Mapping Relational Identity}

\begin{center}
\begin{tikzpicture}
    \fill[AccentCyan] (-2,0) rectangle (-1.7,0.08);
    \fill[AccentCyan!70] (-1.5,0) rectangle (-1.2,0.08);
    \fill[AccentCyan!50] (-1,0) rectangle (-0.7,0.08);
    \fill[AccentCyan!30] (-0.5,0) rectangle (-0.2,0.08);
    \fill[AccentCyan!15] (0,0) rectangle (0.3,0.08);
\end{tikzpicture}
\end{center}

\vspace{0.5cm}

\begin{chapteroverview}
This chapter explores social media platforms---particularly Twitter/X and LinkedIn---as sources for mapping \textbf{relational identity}: how organizations position themselves through their network connections, content sharing, and engagement patterns. Unlike websites (self-presentation) or news (external attribution), social media reveals how actors perform their roles through ongoing digital interactions.
\end{chapteroverview}


%═══════════════════════════════════════════════════════════════════════════════
\section*{Why Social Media?}
%═══════════════════════════════════════════════════════════════════════════════

Social media provides unique insights into role constellations:

\sidenote{\textbf{Social Media Signals:} Following/follower networks, content sharing patterns, engagement interactions, hashtag communities, real-time positioning.}

\begin{description}[leftmargin=0.5cm, itemsep=0.3em]
\item[Network Visibility] Who follows whom, who retweets whom---explicit relational claims
\item[Real-time Dynamics] Continuous stream reveals positioning as it happens
\item[Discourse Participation] Hashtags and mentions show topical alignment
\item[Community Boundaries] Clustering patterns reveal ecosystem structure
\end{description}


%═══════════════════════════════════════════════════════════════════════════════
\section*{Twitter/X as Data Source}
%═══════════════════════════════════════════════════════════════════════════════

Twitter has been the dominant platform for social media research due to its public nature and (historically) accessible API.

\subsection*{Data Types}

\begin{itemize}[leftmargin=*, itemsep=0.2em]
\item \textbf{Follower networks}: Directed ties showing attention flows
\item \textbf{Retweet networks}: Endorsement and amplification
\item \textbf{Mention networks}: Direct engagement
\item \textbf{Content streams}: Tweets revealing discourse positioning
\item \textbf{Hashtag co-occurrence}: Topic communities
\end{itemize}

\begin{criticalbox}[API Access Changes]
Twitter's API access policies changed dramatically in 2023, restricting academic research access. Historical data remains valuable, but ongoing collection faces new barriers.
\end{criticalbox}


%═══════════════════════════════════════════════════════════════════════════════
\section*{LinkedIn for Professional Networks}
%═══════════════════════════════════════════════════════════════════════════════

LinkedIn offers complementary data on professional relationships:

\begin{itemize}[leftmargin=*, itemsep=0.2em]
\item Employee flows between organizations
\item Advisory relationships
\item Institutional affiliations
\item Professional endorsements
\end{itemize}

\sidenote{\textbf{Access Limitation:} LinkedIn's terms of service restrict automated data collection. Most research relies on manual collection or partnership access.}


%═══════════════════════════════════════════════════════════════════════════════
\section*{Analytical Approaches}
%═══════════════════════════════════════════════════════════════════════════════

\begin{description}[leftmargin=0.5cm, itemsep=0.3em]
\item[Network Analysis] Centrality, community detection, brokerage identification
\item[Content Analysis] Topic modeling on posts, sentiment tracking
\item[Temporal Dynamics] How networks and discourse evolve
\item[Bot Detection] Filtering automated accounts for authentic signals
\end{description}


%═══════════════════════════════════════════════════════════════════════════════
\section*{Chapter Summary}
%═══════════════════════════════════════════════════════════════════════════════

\begin{summarybox}
\begin{itemize}[leftmargin=*, itemsep=0.3em]
\item \textbf{Social media} reveals relational identity through ongoing digital interactions.

\item \textbf{Twitter/X} provided rich network and content data, though API access has become restricted.

\item \textbf{LinkedIn} offers professional network data with access limitations.

\item \textbf{Analytical value} lies in real-time dynamics, community detection, and discourse positioning.

\item \textbf{Integration} with other modalities validates claimed relationships and reveals gaps between assertion and engagement.
\end{itemize}
\end{summarybox}

\vspace{0.5cm}

\begin{center}
\begin{tikzpicture}
    \fill[AccentCyan] (-2,0) rectangle (-1.7,0.08);
    \fill[AccentCyan!70] (-1.5,0) rectangle (-1.2,0.08);
    \fill[AccentCyan!50] (-1,0) rectangle (-0.7,0.08);
    \fill[AccentCyan!30] (-0.5,0) rectangle (-0.2,0.08);
    \fill[AccentCyan!15] (0,0) rectangle (0.3,0.08);
\end{tikzpicture}
\end{center}

\closeintermezzo

\end{document}
